% ============================================================
% CHAPTER: PROBLEM DEFINITION
% ============================================================

\chapter{Problem Definition}

In today's fast-paced and highly competitive environment, mental health challenges such as stress, depression, and emotional instability have become increasingly prevalent among students, professionals, and individuals across all age groups. These conditions negatively affect mental and physical health, leading to reduced productivity, poor decision-making, strained relationships, and increased healthcare costs. The rising burden of mental health issues has created a significant need for efficient and accessible early detection systems.

Existing detection methods, including self-reported questionnaires, psychological assessments, and clinical evaluations, often suffer from limitations such as subjectivity, high cost, delayed diagnosis, and lack of continuous monitoring. Many individuals also hesitate to seek professional help due to social stigma and limited accessibility to mental healthcare services. As a result, early signs of stress and depression frequently go unnoticed, leading to severe long-term consequences.

This creates a major gap in the healthcare and wellness sector, highlighting the need for innovative, affordable, and non-invasive detection solutions. Vocal parameter analysis has emerged as a promising approach, as human speech contains valuable indicators of emotional and psychological states through variations in pitch, tone, energy, and speech patterns.

This project aims to address this problem by developing an AI-based multimodal stress detection system using vocal and physical parameters for early identification of stress, depression, and emotional imbalance. The proposed solution offers strong potential not only in improving mental healthcare but also in creating entrepreneurial opportunities in digital health, wellness technology, and AI-driven healthcare solutions.
